% Generated from validated Article Draft Result. Do not hand-edit; revise the structured draft instead.
\section{モデルへの入口が増える――API、open weights、独立プロバイダ}
\noindent\textbf{2021年には、モデルを『使える』ことの意味が一つではなくなった。OpenAI APIのavailability、GPT-J-6Bのopen-weight release、Jurassic-1 / AI21 Studioの独立APIは、それぞれ異なるaccess surfaceである。model capabilityとdistribution modelを分けて見ることで、生成AIが研究成果から利用可能な基盤へ移る過程が見えてくる。}\autocite{src-1b063a3564f3,src-b50beb89f16f,src-bf581c89a17f}

OpenAI APIのno-waitlist availabilityは、model researchとは別のaccess boundaryとして重要である。モデルの論文や能力評価が存在することと、外部の利用者がservice interfaceを通じてアクセスできることは同じ出来事ではない。2021年のdistributionを読む際には、このAPI availabilityをmodel releaseとは分離し、利用可能性そのものが生成AIの普及を左右する独立した層になったと捉える必要がある。\autocite{src-b50beb89f16f}


GPT-J-6Bは、open-weight releaseという別の入口を記録する。API越しにproviderの実行環境を使う形とは異なり、weightが配布されることは利用者側でmodel artifactへ接近できる範囲を変える。ただし、weightが取得可能であることだけを理由に『open source』と同義にしてはいけない。source code、license、training data、再現性などは別の論点であり、本稿ではopen-weightという確認できる境界に留める。\autocite{src-1b063a3564f3}


Jurassic-1 / AI21 Studioは、frontier-scale language modelへのAPI accessが一社だけのsurfaceではなくなり、独立providerからも提供される方向を示す補助線になる。ここで比較したいのはモデル性能の優劣ではなく、利用者から見た入口の構造である。commercial APIという形はopen weightsとは異なり、model artifactそのものを配布せずに能力をserviceとして提供するdistribution modelである。\autocite{src-bf581c89a17f}


この三者を並べると、2021年には少なくとも『provider-managed APIへのavailability』『open-weight artifactへのaccess』『別providerによるAPI service』という異なるdistribution surfaceが同時に見えていたと整理できる。重要なのは、open/closedの二分法へ急いで落とし込まないことである。誰が実行環境を持つか、weightへ触れられるか、serviceとして提供されるかは別々の軸であり、modelの技術的能力とも切り分けて評価すべきである。\autocite{src-1b063a3564f3,src-b50beb89f16f,src-bf581c89a17f}


\begin{claimboundary}
本稿でいうopen weightは、weightへのaccessが確認できるという限定された意味で用いる。open source、完全な再現可能性、training dataの公開、自由な再配布までを自動的に含めない。また各providerが示す能力評価は独立検証済みの事実へ変換しない。\autocite{src-1b063a3564f3,src-b50beb89f16f,src-bf581c89a17f}
\end{claimboundary}
